\section{Related Work}

\paragraph{Language-Guided Web Agent Benchmarks}
The development of reproducible environments for autonomous agents has seen considerable progress in recent years. Earlier efforts introduced reinforcement learning environments~\citep{brockman2016openai}, and extended into web domains~\citep{shi2017world,liu2018reinforcement}. %, which laid the groundwork for reproducible web-based environments. 
Recent web agent benchmarks introduced tasks involving actions on static internet pages~\citep{deng2023mind2web} as well as interaction in simulated web environments~\citep{yao2022webshop, zhou2023webarena}. AgentBench~\citep{liu2023agentbench} extends the scope of agents for computer interaction beyond the web, exploring database management and operating system functionalities. 


\paragraph{LLM Agents}
There has been significant recent interest in using Large Language Models (LLMs) for developing autonomous agents~\citep{xi2023rise,wang2023survey}. State-of-the-art LLMs~\citep{team2023gemini,openai2023gpt4,chowdhery2023palm,rae2021gopher,zhang2022opt,touvron2023llama,touvron2023llama2,jiang2023mistral,jiang2024mixtral} based on the Transformer~\citep{vaswani2017attention} architecture have demonstrated impressive abilities in learning from in-context examples~\citep{brown2020language,chan2022transformers}, reasoning~\citep{wei2022chain,yao2023tree,wang2022self,besta2023graph}, following instructions~\citep{chung2022scaling,longpre2023flan,ouyang2022training}, and operating over long-context sequences~\citep{tay2020long,bertsch2023unlimiformer,tworkowski2023focused}. Several recent works leverage these abilities for building autonomous web agents: \citet{kim2023language} propose a recursive prompting method to improve GPT-4 performance on MiniWoB++~\citep{liu2018reinforcement}. \citet{liu2023bolaa} propose a method of orchestrating multiple LLM agents to improve performance on the WebShop~\citep{yao2022webshop} environment. \citet{zeng2023agenttuning} fine-tunes the LLaMA-2 models on interaction trajectories with instructions, improving over baseline agents.

\paragraph{Vision-Language Models}
Finally, our work builds off advances in vision-language models (VLMs), used for many multimodal tasks such as image captioning~\citep{vinyals2015show}, visual question answering~\citep{antol2015vqa}, and other benchmarks~\citep{mialon2023gaia,yue2023mmmu,tong2024eyes}.
% There have been significant efforts in the modeling space to develop stronger VLMs for solving general vision-language problems. 
Frozen~\citep{tsimpoukelli2021multimodal} was one of the first approaches to demonstrate the effectiveness of finetuning a visual encoder to map images into the embedding space of a LLM, introducing compelling few-shot multimodal abilities. \citet{alayrac2022flamingo} introduced cross-attention layers and scaled up model sizes and training data. \citet{wang2023cogvlm} introduced trainable visual expert modules to improve vision-language fusion. \citet{liu2023visual} proposed fine-tuning on images paired with instructions to improve text generation performance on several multimodal tasks. GPT-4V~\citep{openai2023gpt4} introduces visual processing to the GPT-4 family of models. 
% , unlocking many compelling abilities~\citep{yang2023dawn,yang2023som}. 
Gemini~\citep{team2023gemini} is multimodal from the beginning (in contrast to post-hoc fine-tuned models), and can handle text interleaved with visual and audio inputs. 
Several recent work have also explored using VLMs to build agents for mobile platforms~\citep{zhan2023you,chu2023mobilevlm,yang2023appagent} and the web~\citep{gur2023real,hong2023cogagent}. \citet{zheng2024gpt} is contemporaneous work which performs action grounding to identify appropriate HTML elements for enabling agents to execute actions. In contrast, our proposed SoM agent uses JavaScript to produce a Set-of-Marks~\citep{yang2023som} for the VLM to directly use as an observation and action space.
